\documentclass[10pt,a4paper]{amsart}
\usepackage[margin=19mm]{geometry}
\usepackage{amsmath,amssymb,mathtools}
\usepackage[scr=rsfs,cal=euler]{mathalfa}
\usepackage{xfrac}
\usepackage[unicode,hidelinks,bookmarks=false]{hyperref}
\newcommand{\ie}{i.e.\ }
\newcommand{\cf}{cf.\ }
\newcommand{\resp}{respectively\ }
\newcommand{\Subsection}{\subsection}
\providecommand{\nobreakdash}{\nobreak\hbox{-}}
\setlength{\emergencystretch}{2em}
\multlinegap0pt
\usepackage{amssymb}
\newtheorem{theorem}{Theorem}
\newtheorem{lemma}{Lemma}
\title{Concentrated solutions to fractional Schr\"{o}dinger-Poisson system with non-homogeneous potentials}
\author{Lintao Liu, Haidong Yang}
\date{Source: arXiv:2603.14490v1 (CC BY 4.0)}
\begin{document}
\maketitle

\noindent\emph{Source and licence.} Concentrated solutions to fractional Schr\"{o}dinger-Poisson system with non-homogeneous potentials. Version arXiv:2603.14490v1 (CC BY 4.0), distributed by arXiv under the Creative Commons Attribution 4.0 International licence, http://creativecommons.org/licenses/by/4.0/, as recorded in the arXiv licence metadata for this exact version and shown on the abstract page https://arxiv.org/abs/2603.14490v1. Full licence notice: \texttt{licenses/arxiv-2603.14490-CC-BY-4.0.txt}.\par\smallskip

\noindent\textit{Editorial scope.} Counted unit: the article's lemma lem4-4 (source0.tex lines 1284--1286). The article's complete original proof of it is delivered with it (source0.tex lines 1287--1396, one \texttt{proof} environment of 110 lines). No mathematics is written, repaired or re-derived: the statement and the proof are the article's own lines, in the article's own notation, and no retained line is substituted or re-flowed.  Retained uncounted as the source-local context the proof needs: lines 308--310; lines 340--342; lines 1043--1045; lines 1051--1057; lines 1060--1062; lines 1136--1138; lines 1210--1212; lines 1217--1220; lines 1441--1443; lines 1760--1765.  Nothing was omitted from any retained span, so the package carries no omission record.  No retained line contains a citation, so the wrapper carries no bibliography and no bibliography entry.  No retained line contains a figure environment, an image-inclusion command or drawing code, so no image is delivered and the package declares no asset dependency.  Editorial formatting only, in addition: an AMS \texttt{amsart} preamble that restates the article's own declarations and macros used by the retained lines, namely the environments \texttt{theorem}, \texttt{lemma} on separate counters, together with the standard packages those lines need.  The retained environments print in this document as Theorem 1, Lemma 1 in order of appearance, whereas the article prints the same items with the numbering of its own full document.  The complete native source of the article as distributed in the arXiv e-print for this exact version is bundled unchanged as \texttt{sources/196403/source0.tex}.
	\begin{equation}\label{equ13}
		H(y):=	\int_{\mathbb{R}^{3}}V_{0}(x+y)Q^{2}(x)dx\quad \mathrm {where} \quad V_{0}(x):=V_{i}(x) \quad \mathrm {and }~~i~~\mathrm {satisfying }~~ \bar{\lambda}_{i}=\bar{\lambda}_{0},
	\end{equation}

	\begin{theorem}\label{the3}
		Assume that $V(x)$ satisfies $\mathbf{(V_1)}$-$\mathbf{(V_3)}$, then there exists a unique nonnegative minimizer of $e_{1}(a_{k})$ as $k\rightarrow\infty$.
	\end{theorem}

	\begin{equation}\label{equ4-4}
		\hat{u}_{ik}(x):=\sqrt{a^{\ast}}\varepsilon^{\frac{3}{2}}_{k}u_{ik}(\varepsilon_{k}x+x_{1k})\quad \hbox{in $\mathbb{R}^{3}$,\quad $i=1, 2$}
	\end{equation}

	\begin{lemma}\label{lem4-1}
		Under the assumptions of Theorem \ref{the3}, then there exists a constant $C>0$, independent
		of $k>0$, such that
		\begin{equation}\label{equ4-6}
			\int_{\mathbb{R}^{3}}|(-\Delta)^{\frac{s}{2}}\hat{\eta}_{k}|^{2}dx+\varepsilon_{k}^{2s}\int_{\mathbb{R}^{3}}V(\varepsilon_{k}x+x_{1k})\hat{\eta}_{k}^{2}dx+\int_{\mathbb{R}^{3}}\hat{\eta}_{k}^{2}dx\leq C\quad \hbox{as $k\rightarrow\infty$}.
		\end{equation}
	\end{lemma}

		\begin{equation}\label{equ4-7}
			(-\Delta)^{s} \hat{u}_{ik}+\varepsilon_{k}^{2s}V(\varepsilon_{k}x+x_{1k})\hat{u}_{ik}=\mu_{ik}\varepsilon_{k}^{2s}\hat{u}_{ik}+\frac{\varepsilon_{k}^{4s-3}}{a^{\ast}}\int_{\mathbb{R}^{3}}\frac{\hat{u}_{ik}^{2}(y)}{|x-y|^{3-2s}}dy\hat{u}_{ik}+\hat{u}_{ik}^{p-1}\quad \hbox{in $\mathbb{R}^{3}$},
		\end{equation}

		\begin{equation}\label{equ4-16}
			|\hat{\eta}_{k}(x)|\leq\frac{C}{1+|x|^{3+2s}}\quad \hbox{as $k\rightarrow\infty$}.
		\end{equation}

	\begin{equation}\label{equ4-22}
		\hat{\eta}_{0}(x)=b_{0}\left(Q+\frac{p-2}{2s}x\cdot\nabla Q\right)+\sum_{i=1}^{3}b_{i}\frac{\partial Q}{\partial x_{i}},
	\end{equation}

		\begin{equation}\label{equ4-23}
			\frac{b_{0}(p-2)}{2s}\int_{\mathbb{R}^{3}}\frac{\partial V_{0}(x+y_{0})}{\partial x_{j}}(x\cdot\nabla Q^{2})dx-\sum^{3}_{i=1}
			b_{i}\int_{\mathbb{R}^{3}}\frac{\partial^{2}V_{0}(x+y_{0})}{\partial x_{i}\partial x_{j}}Q^{2}dx=0.
		\end{equation}

		\begin{equation}\label{equ5-1}
			\hat{u}_{ik}(x)+	|\nabla\hat{u}_{ik}(x)|\leq\frac{C}{1+|x|^{3+2s}}\quad \hbox{as $k\rightarrow\infty$}.
		\end{equation}

	\begin{lemma}\label{lem5-4}
		Assume that $x_{1k}$ is the unique global maximum point of $u_{1k}$, and $\hat{u}_{ik}$ is given in \eqref{equ4-4}, then there holds that
		\begin{equation*}
			\int_{\mathbb{R}^{3}}(-\Delta)^{s}\hat{u}_{ik}[(x-x_{1k})\cdot\nabla\hat{u}_{ik}]dx=\frac{2s-3}{2} \int_{\mathbb{R}^3} \left| (-\Delta)^{\frac{s}{2}}\hat{u}_{ik}\right| ^2dx.
		\end{equation*}
	\end{lemma}

	\begin{lemma}\label{lem4-4}
		Under the assumptions of Theorem \ref{the3}, and assume that $V_{0}(x)$ is defined by \eqref{equ13}, then we have that $\hat{\eta}_{0}\equiv0$.
	\end{lemma}

	\begin{proof}
		Multiplying \eqref{equ4-7} by $(x-x_{1k})\cdot\nabla\hat{u}_{ik}(x)$ and integrating over $\mathbb{R}^{3}$, it follows that
		\begin{equation}\label{equ4-29}
			\begin{split}
				&\int_{\mathbb{R}^{3}}(-\Delta)^{s}\hat{u}_{ik}[(x-x_{1k})\cdot\nabla\hat{u}_{ik}]dx+\varepsilon_{k}^{2s}\int_{\mathbb{R}^{3}}V(\varepsilon_{k}x+x_{1k})\hat{u}_{ik}[(x-x_{1k})\cdot\nabla\hat{u}_{ik}]dx\\
				=&\mu_{ik}\varepsilon_{k}^{2s}\int_{\mathbb{R}^{3}}\hat{u}_{ik}[(x-x_{1k})\cdot\nabla\hat{u}_{ik}]dx+\int_{\mathbb{R}^{3}}\hat{u}_{ik}^{p-1}[(x-x_{1k})\cdot\nabla\hat{u}_{ik}]dx\\
				&+\frac{\varepsilon_{k}^{4s-3}}{a^{\ast}}\int_{\mathbb{R}^{3}}\int_{\mathbb{R}^{3}}\frac{\hat{u}_{ik}^{2}(y)}{|x-y|^{3-2s}}\hat{u}_{ik}(x)[(x-x_{1k})\cdot\nabla\hat{u}_{ik}]dydx.
			\end{split}
		\end{equation}
		Applying \eqref{equ4-7} and Lemma \ref{lem5-4}, one deduce that
		\begin{equation}\label{equ4-30}
			\begin{split}
				&\int_{\mathbb{R}^{3}}(-\Delta)^{s}\hat{u}_{ik}[(x-x_{1k})\cdot\nabla\hat{u}_{ik}]dx=\frac{2s-3}{2} \int_{\mathbb{R}^3} \left| (-\Delta)^{\frac{s}{2}}\hat{u}_{ik}\right| ^2dx\\
				=&-\frac{2s-3}{2}\varepsilon_{k}^{2s}\int_{\mathbb{R}^{3}} V(\varepsilon_{k}x+x_{1k})\hat{u}_{ik}^{2}dx+ \frac{2s-3}{2}\mu_{ik}\varepsilon_{k}^{2s}\int_{\mathbb{R}^{3}}\hat{u}_{ik}^{2}dx\\
				&+\frac{2s-3}{2}\int_{\mathbb{R}^{3}}\hat{u}_{ik}^{p}dx+\frac{2s-3}{2}\frac{\varepsilon_{k}^{4s-3}}{a^{\ast}}\int_{\mathbb{R}^{3}}\int_{\mathbb{R}^{3}}\frac{\hat{u}_{ik}^{2}(y)\hat{u}_{ik}^{2}(x)}{|x-y|^{3-2s}}dydx.
			\end{split}
		\end{equation}
		Integrating by parts, one can see that
		\begin{equation}\label{equ4-31}
			\begin{split}
				&\varepsilon_{k}^{2s}\int_{\mathbb{R}^{3}}V(\varepsilon_{k}x+x_{1k})\hat{u}_{ik}[(x-x_{1k})\cdot\nabla\hat{u}_{ik}]dx\\
				=&\varepsilon_{k}^{2s}\lim_{R\rightarrow\infty}\int_{ B_{R}(0)}V(\varepsilon_{k}x+x_{1k})\hat{u}_{ik}[(x-x_{1k})\cdot\nabla\hat{u}_{ik}]dx\\
				=&\frac{\varepsilon_{k}^{2s}}{2}\lim_{R\rightarrow\infty}\int_{\partial B_{R}(0)}V(\varepsilon_{k}x+x_{1k})(x-x_{1k})\cdot\nu\hat{u}_{ik}^{2}dS\\
				&-\frac{\varepsilon_{k}^{2s}}{2}\int_{\mathbb{R}^{3}}\left[3V(\varepsilon_{k}x+x_{1k})+(x-x_{1k})\cdot\nabla V(\varepsilon_{k}x+x_{1k})\right] \hat{u}_{ik}^{2}dx\\
				=&-\frac{\varepsilon_{k}^{2s}}{2}\int_{\mathbb{R}^{3}}\left[3V(\varepsilon_{k}x+x_{1k})+(x-x_{1k})\cdot\nabla V(\varepsilon_{k}x+x_{1k})\right] \hat{u}_{ik}^{2}dx,
			\end{split}
		\end{equation}
		\begin{equation}\label{equ4-32}
			\begin{split}
				&\mu_{ik}\varepsilon_{k}^{2s}\int_{\mathbb{R}^{3}}\hat{u}_{ik}[(x-x_{1k})\cdot\nabla\hat{u}_{ik}]dx=\mu_{ik}\varepsilon_{k}^{2s}\lim_{R\rightarrow\infty}\int_{ B_{R}(0)}\hat{u}_{ik}[(x-x_{1k})\cdot\nabla\hat{u}_{ik}]dx\\
				=&\frac{\mu_{ik}\varepsilon_{k}^{2s}}{2}\lim_{R\rightarrow\infty}\int_{\partial B_{R}(0)}(x-x_{1k})\cdot\nu\hat{u}_{ik}^{2}dS-\frac{3\mu_{ik}\varepsilon_{k}^{2s}}{2}\int_{\mathbb{R}^{3}}\hat{u}_{ik}^{2}dx=-\frac{3\mu_{ik}\varepsilon_{k}^{2s}}{2}\int_{\mathbb{R}^{3}}\hat{u}_{ik}^{2}dx,
			\end{split}
		\end{equation}
		\begin{equation}\label{equ4-33}
			\begin{split}
				&\int_{\mathbb{R}^{3}}\hat{u}_{ik}^{p-1}[(x-x_{1k})\cdot\nabla\hat{u}_{ik}]dx=\lim_{R\rightarrow\infty}\int_{ B_{R}(0)}\hat{u}_{ik}^{p-1}[(x-x_{1k})\cdot\nabla\hat{u}_{ik}]dx\\
				=&\frac{1}{p}\lim_{R\rightarrow\infty}\int_{\partial B_{R}(0)}(x-x_{1k})\cdot\nu\hat{u}_{ik}^{p}dS-\frac{3}{p}\int_{\mathbb{R}^{3}}\hat{u}_{ik}^{p}dx=-\frac{3}{p}\int_{\mathbb{R}^{3}}\hat{u}_{ik}^{p}dx,
			\end{split}
		\end{equation}
		and
		\begin{equation}\label{equ4-34}
			\begin{split}
				&\frac{\varepsilon_{k}^{4s-3}}{a^{\ast}}\int_{\mathbb{R}^{3}}\int_{\mathbb{R}^{3}}\frac{\hat{u}_{ik}^{2}(y)}{|x-y|^{3-2s}}\hat{u}_{ik}(x)[(x-x_{1k})\cdot\nabla\hat{u}_{ik}]dydx\\
				=&\frac{\varepsilon_{k}^{4s-3}}{a^{\ast}}\lim_{R\rightarrow\infty}\int_{ B_{R}(0)}\int_{\mathbb{R}^{3}}\frac{\hat{u}_{ik}^{2}(y)}{|x-y|^{3-2s}}\hat{u}_{ik}(x)[(x-x_{1k})\cdot\nabla\hat{u}_{ik}]dydx\\
				=&\frac{\varepsilon_{k}^{4s-3}}{2a^{\ast}}\lim_{ R\rightarrow\infty}\int_{ \partial B_{R}(0)}\int_{\mathbb{R}^{3}}\frac{\hat{u}_{ik}^{2}(y)}{|x-y|^{3-2s}}(x-x_{1k})\cdot\nu\hat{u}_{ik}^{2}(x)dydS\\
				&-\frac{3\varepsilon_{k}^{4s-3}}{2a^{\ast}}\int_{\mathbb{R}^{3}}\int_{\mathbb{R}^{3}}\frac{\hat{u}_{ik}^{2}(y)\hat{u}_{ik}^{2}(x)}{|x-y|^{3-2s}}dydx\\
				&+\frac{(3-2s)\varepsilon_{k}^{4s-3}}{2a^{\ast}}\int_{\mathbb{R}^{3}}\int_{\mathbb{R}^{3}}\frac{(x-x_{1k})\cdot(x-y)}{|x-y|^{5-2s}}\hat{u}^{2}_{ik}(x)\hat{u}^{2}_{ik}(y)dydx\\
				=&-\frac{(3+2s)\varepsilon_{k}^{4s-3}}{4a^{\ast}}\int_{\mathbb{R}^{3}}\int_{\mathbb{R}^{3}}\frac{\hat{u}_{ik}^{2}(y)\hat{u}_{ik}^{2}(x)}{|x-y|^{3-2s}}dydx,
			\end{split}
		\end{equation}
		where \eqref{equ4-31}-\eqref{equ4-34} hold due to the facts that
		\begin{equation*}
			\left| \int_{\partial B_{R}(0)}V(\varepsilon_{k}x+x_{1k})(x-x_{1k})\cdot\nu\hat{u}_{ik}^{2}dS\right|
			\leq CR^{-6-4s}R^{2}\rightarrow0\quad \hbox{as $R\rightarrow\infty$},
		\end{equation*}
		\begin{equation*}
			\left| \int_{\partial B_{R}(0)}(x-x_{1k})\cdot\nu\hat{u}_{ik}^{2}dS\right|
			\leq CR^{-6-4s}R^{2}\rightarrow0\quad \hbox{as $R\rightarrow\infty$},
		\end{equation*}
		\begin{equation*}
			\left| \int_{\partial B_{R}(0)}(x-x_{1k})\cdot\nu\hat{u}_{ik}^{p}dS\right| \leq CR^{-p(3+2s)}R^{2}\rightarrow0\quad \hbox{as $R\rightarrow\infty$},
		\end{equation*}
		\begin{align*} 
			&	\left| \int_{ \partial B_{R}(0)}\int_{\mathbb{R}^{3}}\frac{\hat{u}_{ik}^{2}(y)}{|x-y|^{3-2s}}(x-x_{1k})\cdot\nu\hat{u}_{ik}^{2}(x)dydS\right| \\
			\leq&\left| \int_{\mathbb{R}^{3}}\frac{\hat{u}_{ik}^{2}(y)}{|x-y|^{3-2s}}dy\int_{ \partial B_{R}(0)}(x-x_{1k})\cdot\nu\hat{u}_{ik}^{2}(x)dS\right| \leq CR^{-6-4s}R^{2}\rightarrow0\quad \hbox{as $R\rightarrow\infty$},
		\end{align*}
		and
		\begin{align*} 
			\int_{\mathbb{R}^{3}}\int_{\mathbb{R}^{3}}\frac{(x-x_{1k})\cdot(x-y)}{|x-y|^{5-2s}}\hat{u}^{2}_{ik}(x)\hat{u}^{2}_{ik}(y)dydx
			=\frac{1}{2}\int_{\mathbb{R}^{3}}\int_{\mathbb{R}^{3}}\frac{\hat{u}^{2}_{ik}(x)\hat{u}^{2}_{ik}(y)}{|x-y|^{3-2s}}dydx.
		\end{align*}
		It then follows from \eqref{equ4-29}-\eqref{equ4-34} that
		\begin{align*} 
			\left( \frac{3}{p}+\frac{2s-3}{2}\right) \int_{\mathbb{R}^{3}}\hat{u}_{ik}^{p}dx=&sa^{\ast}+s\varepsilon_{k}^{2s}\int_{\mathbb{R}^{3}}V(\varepsilon_{k}x+x_{1k})\hat{u}_{ik}^{2}dx\\
			&+\frac{1}{2}\varepsilon_{k}^{2s}\int_{\mathbb{R}^{3}}\left[ (x-x_{1k})\cdot \nabla V(\varepsilon_{k}x+x_{1k})\right] \hat{u}_{ik}^{2}dx\\
			&-\frac{(6s-3)\varepsilon_{k}^{4s-3}}{4a^{\ast}}\int_{\mathbb{R}^{3}}\int_{\mathbb{R}^{3}}\frac{\hat{u}_{ik}^{2}(y)\hat{u}_{ik}^{2}(x)}{|x-y|^{3-2s}}dydx \quad \hbox{as $k\rightarrow\infty$},
		\end{align*}
		where we have used the facts that $\mu_{ik}\varepsilon^{2s}_{k}\rightarrow-1$ and $\hat{u}_{ik}(x)\rightarrow Q(x)$ in $L^{\infty}(\mathbb{R}^{3})$
		as $k\rightarrow\infty$. Thus we have
		\begin{equation*}
			\begin{split}
				&\frac{6+(2s-3)p}{2}\int_{\mathbb{R}^{3}}\left[\theta \hat{u}_{1k}+(1-\theta)\hat{u}_{2k}\right]^{p-1}\hat{\eta}_{k}dx\\
				=&s\varepsilon_{k}^{2s}\int_{\mathbb{R}^{3}}V(\varepsilon_{k}x+x_{1k})(\hat{u}_{1k}+\hat{u}_{2k})\hat{\eta}_{k}dx\\
				&+\frac{1}{2}\varepsilon_{k}^{2s}\int_{\mathbb{R}^{3}}\left[ (x-x_{1k})\cdot \nabla V(\varepsilon_{k}x+x_{1k})\right](\hat{u}_{1k}+\hat{u}_{2k})\hat{\eta}_{k}dx\\
				&-\frac{(6s-3)\varepsilon_{k}^{4s-3}}{4a^{\ast}}\int_{\mathbb{R}^{3}}\int_{\mathbb{R}^{3}}\frac{\hat{u}_{1k}^{2}(x)}{|x-y|^{3-2s}}(\hat{u}_{1k}(y)+\hat{u}_{2k}(y))\hat{\eta}_{k}(y)dydx\\
				&-\frac{(6s-3)\varepsilon_{k}^{4s-3}}{4a^{\ast}}\int_{\mathbb{R}^{3}}\int_{\mathbb{R}^{3}}\frac{\hat{u}_{2k}^{2}(y)}{|x-y|^{3-2s}}(\hat{u}_{1k}(x)+\hat{u}_{2k}(x))\hat{\eta}_{k}(x)dydx:=A_{k}(x).
			\end{split}
		\end{equation*}
		Using polynomial decay \eqref{equ4-16} and \eqref{equ5-1}, similar arguments of Lemma \ref{lem4-1} yield that
		\begin{equation*}
			A_{k}(x)=O(\varepsilon_{k}^{4s-3})\quad \hbox{as $k\rightarrow\infty$}.
		\end{equation*}
		By the fact that $\hat{u}_{ik}(x)\rightarrow Q(x)$ in $L^{\infty}(\mathbb{R}^{3})$ as $k\rightarrow\infty$, we then deduce from \eqref{equ4-16}, \eqref{equ4-22} and \eqref{equ5-1} that
		\begin{equation*}
			\begin{split}
				0&=\int_{\mathbb{R}^{3}}Q^{p-1}\hat{\eta}_{0}dx\\
				&=\int_{\mathbb{R}^{3}}Q^{p-1}\left[ b_{0}\left(Q+\frac{p-2}{2s}x\cdot\nabla Q\right)+\sum_{i=1}^{3}b_{i}\frac{\partial Q}{\partial x_{i}}\right] dx\\
				&=b_{0}\int_{\mathbb{R}^{3}}Q^{p}dx+\frac{b_{0}(p-2)}{2sp}\int_{\mathbb{R}^{3}}x\cdot\nabla Q^{p}dx+\sum_{i=1}^{3}\frac{b_{i}}{p}\int_{\mathbb{R}^{3}}\frac{\partial Q^{p}}{\partial x_{i}}dx\\
				&=b_{0}\int_{\mathbb{R}^{3}}Q^{p}dx-\frac{3b_{0}(p-2)}{2sp}\int_{\mathbb{R}^{3}}Q^{p}dx\\
				&=\frac{2sp-3p+6}{2sp}b_{0}\int_{\mathbb{R}^{3}}Q^{p}dx,
			\end{split}
		\end{equation*}
		which implies that $b_{0}=0$ due to $2<p<\frac{4s}{3}+2<\frac{6}{3-2s}$. Substituting $b_{0}=0$ into \eqref{equ4-23}, one derive that
		\begin{equation*}
			\sum^{3}_{i=1}
			b_{i}\int_{\mathbb{R}^{3}}\frac{\partial^{2}V_{0}(x+y_{0})}{\partial x_{i}\partial x_{j}}Q^{2}dx=0.
		\end{equation*}
		Since $y_{0}$ is the unique and non-degenerate critical point of
		$H(y)$, then we can obtain that $b_{1}=b_{2}=b_{3}=0$, and Lemma \ref{lem4-4} is proved.
	\end{proof}

\end{document}
